% ABOUTME: Deep learning test set results and the ML vs DL comparison of the IEEE report (Phase 2).
% ABOUTME: Frozen networks and hybrids scored once on FD001 test engines over 5 seeds, with bootstrap.
\section{Deep Learning Test Results}
\label{sec:dl-results}

\subsection{Protocol}
Every architecture and hyperparameter, and the choice of the Transformer as
the selected network, were fixed by cross validation
(Sections~\ref{sec:dl-models} and~\ref{sec:hybrid}) before the test set was
opened. Each network is retrained with the same procedure as inside a
cross validation fold, the same 80/20 split into fitting and early stopping
engines, now of all 100 training engines: 80 to fit, 20 held out for early
stopping, scalers fitted on the 80 only, and the weights of the best early
stopping epoch kept. No fixed
epoch count is introduced, so nothing is chosen at this stage. Because a
network's result depends on its initial weights and batch order, every
network is trained with 5 seeds (42 to 46), and the seed also chooses the 20
early stopping engines. We report the mean and standard deviation over the
seeds. Each network predicts once per test engine from the 30 cycle window
ending at its last observed cycle; the hybrids also read the 142 engineered
features at that cycle. As in Section~\ref{sec:results}, predictions are
scored against the true, uncapped RUL. The Phase 1 rows reuse the per engine
predictions of Section~\ref{sec:results} unchanged. All numbers are from
\texttt{results/dl\_test\_metrics.csv} and
\texttt{results/dl\_test\_bootstrap.csv}, produced by
\texttt{notebooks/09\_dl\_results.ipynb}.

\begin{table*}[t]
\centering
\caption{Test results on the 100 FD001 test engines, scored against the
true RUL. Networks: mean $\pm$ standard deviation over 5 seeds. Phase 1 rows
are single deterministic fits on engineered features. Late stage is true
RUL $\le 50$ (33 engines). Capped RMSE scores against test labels clipped at
125, as many published results do. CV RMSE is from cross validation and is on
a different scale (every cycle, capped labels).}
\label{tab:dl-test}
\footnotesize
\setlength{\tabcolsep}{4pt}
\begin{tabular}{lrrrrrrr}
\hline
Model & RMSE & MAE & Score $S$ & Early RMSE & Late RMSE & Capped RMSE & CV RMSE \\
\hline
XGBoost (Phase 1, selected)  & 15.36 & 11.26 & 389 & 18.11 & 7.03 & 14.08 & 13.99 \\
Random Forest (Phase 1)      & 15.23 & 11.33 & 344 & 17.88 & 7.34 & 14.02 & 14.42 \\
\hline
LSTM                         & 16.02 $\pm$ 0.67 & 11.80 & 622 $\pm$ 145 & 19.07 & 6.21 $\pm$ 1.15 & 15.00 & 14.02 \\
Temporal CNN                 & 16.84 $\pm$ 0.93 & 12.51 & 584 $\pm$ 191 & 19.73 & 8.10 $\pm$ 2.12 & 15.64 & 15.91 \\
Transformer (selected)       & 15.37 $\pm$ 0.66 & 11.77 & 422 $\pm$ 98  & 18.30 & 5.88 $\pm$ 1.59 & 14.31 & 13.28 \\
\hline
Hybrid, Transformer + features & 15.92 $\pm$ 0.66 & 12.13 & 450 $\pm$ 65 & 19.00 & 5.85 $\pm$ 0.95 & 14.89 & 13.47 \\
Hybrid, LSTM + features      & 15.93 $\pm$ 0.67 & 11.90 & 493 $\pm$ 178 & 18.97 & 6.11 $\pm$ 1.11 & 14.78 & 14.35 \\
Residual hybrid, XGBoost + Transf. & 15.13 $\pm$ 0.10 & 11.11 & 348 $\pm$ 22 & 18.14 & 5.03 $\pm$ 0.69 & 13.88 & 13.49 \\
\hline
\end{tabular}
\end{table*}

\subsection{Learned sequences against engineered features}
Table~\ref{tab:dl-test} gives the results. The selected Transformer reaches
a test RMSE of $15.37 \pm 0.66$ cycles, the same as the selected Phase 1
model, XGBoost on engineered features (15.36). Cross validation had put the
Transformer 0.71 cycles ahead (13.28 against 13.99); on the test set that
lead is gone. A paired bootstrap over the 100 test engines, in which each
resample is scored with all 5 seeds and their RMSEs averaged, gives a 95\%
interval of $-2.3$ to $+2.2$ cycles for the Transformer minus XGBoost, so the
two are tied. Learning from raw 30 cycle sensor windows therefore matches,
but does not beat, our hand engineered degradation features on FD001. The
LSTM (16.02) and the temporal CNN (16.84) are behind, in the same order as in
cross validation.

The difference between the two approaches is where along the trajectory
the errors fall. Every network except the temporal CNN is more accurate than
XGBoost close to failure, with late stage RMSE between 5.03 and 6.21 against
7.03. For true RUL up to 25 cycles (19 engines), the Transformer's 5 seed
mean prediction has an RMSE of 2.5 cycles against 6.6 for XGBoost. The
networks lose between 51 and 100 cycles, where degradation first becomes
visible: there the Transformer's RMSE is 19 to 22 cycles against 15.7 to
17.5 for XGBoost. Further from failure the Transformer is better again: 12.4
against 16.7 for true RUL from 101 to 125 cycles (22 engines), and 20.0
against 25.0 above the training cap (11 engines). The Transformer also leans late, with a mean error of $+1.9$
cycles and 11 predictions more than 20 cycles late against 9 for XGBoost,
which is why its score $S$ (422) is higher than XGBoost's (389) at the same
RMSE. The late stage holds only 33 engines and was not part of the
comparisons fixed before testing, so the late stage advantage is a
pattern that holds across the seeds rather than a tested result.

\subsection{Hybrids}
Adding the engineered features to the Transformer does not help on the test
set: the two branch hybrid scores $15.92 \pm 0.66$, 0.55 cycles worse than
the Transformer alone, with a bootstrap interval of $-0.2$ to $+1.3$ cycles
that includes zero. Its second place in cross validation did not carry over.
The residual hybrid, XGBoost corrected by a Transformer trained on XGBoost's
residuals, has the lowest test RMSE of all models ($15.13 \pm 0.10$), the
lowest late stage RMSE (5.03) and fewer predictions more than 20 cycles
late than XGBoost (7 against 9). It keeps XGBoost's accuracy in the middle of the trajectory and
most of the Transformer's gain near failure, and it is the most stable model
across seeds, because most of its prediction comes from the deterministic
XGBoost. It was not the model selected by cross validation, so we report it
but do not choose it after the fact. The same bootstrap gives its RMSE minus
XGBoost's as $-0.23$ cycles with an interval of $-1.14$ to $+0.77$; this pair
was not fixed before testing, so the interval is descriptive only, and it
includes zero.

\subsection{Comparison with published results}
Against test labels clipped at 125, the residual hybrid (13.88), our Phase 1
Random Forest (14.02) and XGBoost (14.08) and the Transformer (14.31) rank
seventh to tenth of 21 by RMSE among our models and the published FD001
deep learning results of Table~\ref{tab:dl-published}. They are ahead of
CNN-LSTM (14.40), the attention LSTM with handcrafted features of Chen et
al.\ (14.53) and the five other published models below them, and behind six published deep
models, from BLCNN (13.18) to the GCU-Transformer (11.27). The caveats of
Section~\ref{sec:dl-related} apply: several published values are transcribed
from other papers, and protocols differ in details that move RMSE by about a
cycle.
